\documentclass[10pt]{article}

\usepackage[utf8]{inputenc}

\usepackage[T1]{fontenc}

\usepackage{amsmath}

\usepackage{amsfonts}

\usepackage{amssymb}

\usepackage[version=4]{mhchem}

\usepackage{stmaryrd}

\usepackage{graphicx}

\usepackage[export]{adjustbox}

\graphicspath{ {./images/} }



\begin{document}

имеет место $f_{(r)}(0)=0, r=1,2, \ldots$, но не существует $f_{(1)}(x)=f^{\prime}(x)$ при $x \neq 0$ (ибо $f(x)$ разрывна при $x \neq 0$ ). Следовательно, не существует обычная производная $f^{(r)}(0)$ при $r>1$.



Вводятся также и бесконечные о б о б ш е н в ы $е$ произ во д ны е П е а но. Пусть для всех $t$ с $|t|<\delta$ имеет место



\[

f\left(x_{0}+t\right)=\alpha_{0}+\alpha_{1} t+\ldots+\frac{\alpha_{r-1}}{(r-1) !} t^{r-1}+\frac{\alpha_{r}(t)}{r !} t^{r}

\]



где $\alpha_{0}, \ldots, \alpha_{r-1}-$ постоянные и $\alpha_{r}(t) \rightarrow \alpha_{r}$ при $t \rightarrow 0$ $\left(\alpha_{r}-\right.$ число или символ $\left.\infty\right)$. Тогда $\alpha_{r}$ также наз. П. п. порядка $r$ функции $f$ в точке $x_{0}$. Введена Дж. Пеано (G. Peano).



A. А. Конюиков.



IЕАНО ТЕОРЕМА - одна из теорем существования решения обыкновенного дифференциального уравнения, установленная Дж. Пеано [1] и состоящая в следуюпем. Пусть дано дифференциальное уравнение



\[

y^{\prime}=f(x, y) .

\]



Тогда если функция $t$ ограничена и непрерывна в области $G$, то через каждую внутреннюю точку $\left(x_{0}\right.$, $y_{0}$ ) этой области проходит, шо крайней мере, одна интегральная кривая уравнения (*). Может оказаться, что через нек-рую точку проходит более одной интегральной кривой, напр. для уравнения $y^{\prime}=2 \sqrt{\frac{1}{y}}$ существует бесконечное множество интегральных кривых, проходящих через точку $(0,0)$ :



\[

\begin{gathered}

y(0),-b \leqslant x \leqslant a, \\

y=-(x+b)^{2}, \quad x \leqslant-b, \\

y=(x-a)^{2}, \quad x \geqslant a,

\end{gathered}

\]



где $a, b-$ произвольные постоянные.



Имеются обобщения (в том числе многомерные) П. T. (см. [2], [3]).



Лum.: [1] P e a no G., "Math. Ann.», 1890, Bd 37, S. 182228; [2] П е т р о в с к и й И. Г., Лекции по теории обыкноХ а р т м а н Ф., Обыкновенные дифференциальные уравнения пер. с англ., М., 1970.



ПЕЙДА ТЕОРЕМА - 1) П. т. $\quad 0 \quad$ н $\mathrm{y} \pi$ я $\mathrm{x}$ $L-ф$ у нкц и й Д и рихл е: пусть $L(s, \chi)$-Дирихле L-бункция, $s=\sigma+i t, \chi-$ Дирихле характер по $\bmod d, d \geqslant 3 ;$ существуют аб̆солютные положительные постоянные $c_{1}$, . ., $c_{8}$ такие, что



a) $L(s, \chi) \neq 0$ при $\sigma>1-c_{1} / \log d t, t \geqslant 3$;



б) $L(s, \chi) \neq 0$ при $\sigma>1-c_{2} / \log d, 0<t<5$;



в) для комплексного $\chi \bmod d$,



\[

L(s, \chi) \neq 0 \text { при } \sigma>1-c_{3} / \log d,|t| \leqslant 5 ;

\]



г) для действительного примитивного $\chi \bmod d$,



\[

L(\sigma, \chi) \neq 0 \text { при } \sigma>1-c_{4} / \sqrt{d} \log ^{2} d ;

\]



д) для $2<d \leqslant D$ существует не более одного $d=d_{0}$, $d_{0}>c_{5} \log ^{2} D /(\log \log D)^{8}$ и не более одного действительного примитивного $\chi_{1} \bmod d_{0}$, для к-рого $L\left(s, \chi_{1}\right)$ может иметь действительный нуль $\beta_{1}>1-c_{6} / \log D$, причем $\beta_{1}$ является однократным нулем и для всех $\beta$ таких, что $L(\beta, \chi)=0, \beta>1-c_{6} / \log D$ с действительным $\chi \bmod d$ имеют $d \equiv 0\left(\bmod d_{0}\right)$.



\begin{enumerate}

  \setcounter{enumi}{1}

  \item П. т. о $\pi(x ; d, l)$ - числе простых чисел $p \leqslant x$, $p \equiv l(\bmod d)$, при $0<l \leqslant d, \quad l$ и $d$ - взаимно простых. В обозначениях и при условиях п. 1 вследствие а) в) и д) справедливо равенство

\end{enumerate}



\[

\begin{aligned}

\pi(x ; d, l)= & \frac{\operatorname{li} x}{\varphi(d)}-E \frac{\chi(l)}{\varphi(d)} \sum_{n \leqslant x} \frac{n^{\beta,-1}}{\log n}+ \\

& +O\left(x e^{-c_{7} V \log x}\right),

\end{aligned}

\]



где $E=1$ или $E=0$, смотря шо тому, существует $\beta_{1}$ или нет для данного $d$, и вследствие 2), для любого $d \leqslant(\log x)^{1-\delta}$ при фиксированном $\delta>0$



\[

\pi(x ; d, l)=\frac{\mathrm{li} x}{\varphi(d)}+O\left(x e^{-c_{8} \sqrt{\log x}}\right) .

\]



Этот результат является единственным (1983), к-рый эффективен в том смысле, что, если значение $\delta$ задано, можно указать численные значения $c_{8}$ и постоянной, входящей в символ 0. Замена оценки 2) оценкой Зигеля: $L(\sigma, x) \neq 0$ при $\sigma>1-c(\varepsilon) d^{-\varepsilon}, \varepsilon>0$ распростравяет действие формулы (\textit{) на существенно большие $d, d \leqslant(\log x)^{A}$, с любым фиксированным $A$, но при этом утрачивается эффективность оценок в формуле (}) - по заданному $\varepsilon>0$ невозможно оценить $c_{8}=c_{8}(\varepsilon)$ и $O=O_{\varepsilon}$.



П. т. установлены Э. Пейджем [1].



Jum.: [1] P a g e A., "Proc. London Math. Soc. Ser. 2",



\begin{center}

\includegraphics[max width=\textwidth]{2023_07_08_242f8a9ed2e146ba28cbg-1(1)}

\end{center}



\begin{center}

\includegraphics[max width=\textwidth]{2023_07_08_242f8a9ed2e146ba28cbg-1}

\end{center}



ПЕКЛЕ ЧИСЛО - один из критериев подобия для процессов конвективного теплообмена. П. ч. характеризует соотношение между конвективным и молекулярным процессами переноса тепла в потоке жидкости. П. ч.



\[

\mathrm{Pe}=v l / a=C_{p} \rho v /(\lambda / l),

\]



где $l$ - характерный линейный размер поверхности теплообмена, $v$ - скорость потока жидкости относительно поверхности теплообмена, $a$ - коэффициевт температуропроводности, $C_{p}$ - теплоемкость при постоянном давлении, $\rho$ - плотность и $\lambda$ - коэффициент теплопроводности жидкости.



П. ч. связано с Рейнольдса числом Re II Прандт.ля числом $\operatorname{Pr}$ соотношением $\mathrm{Pe}=\mathrm{Re} \cdot \mathrm{Pr}$.



П. ч. вазвано по имени Ж. Пекле (J. Péclet).



ПЕЛ По материалам одноименной статьи из БCэ-з. ПЕЛлЯ УРАВНЕНИЕ - дюофантово уравнение вида



\[

x^{2}-d y^{2}=1

\]



а также более общее уравнение



\[

x^{2}-d y^{2}=c

\]



где $d$ - натуральное, $\sqrt{d}-$ иррациональное число, $c$ - целое, неизвестные $x$ и $y$ - целые числа.



Если $P_{s} / Q_{s}, s=0,1,2, \ldots,-$ подходящие дроби разложения $\sqrt[V]{d}$ в цепную дробь с периодом $k$, то положительные решения уравнения (1) имеют вид



\[

x=P_{k n-1}, \quad y=Q_{k n-1},

\]



где $n$ - любое натуральное число такое, что $k n$ четно.



Все решения уравнения (1) получаются из формулы



\[

x+y \sqrt{\bar{d}}= \pm\left(x_{0}+y_{0} \sqrt{ } \bar{d}\right)^{n},

\]



где $n$ - любое целое, а $\left(x_{0}, y_{0}\right)$ решение с наименьгшими положительными значениями неизвестных. Общее уравнение (2) либо совсем не имеет решений, либо бесконечно много. При $c=-1$ решения существуют тогда и только тогда, когда $k$ нечетно. При $c=4$ уравнение (2) всегда имеет решения. С помощью решений П. у. при $c= \pm 1, \pm 4$ находятся единицы квадратичвого поля $R(\sqrt{\bar{d}})$. Решения П. у. используются при нахождении автоморфизмов бинарных квадратичных форм $A x^{2}+B x y+C y^{2}$; они позволяют по одному peшению диофантова уравнения $A x^{2}+B x y+C y^{2}=n$ получить бесконечное множество решений.



Уравнение (1) изучалось У. Броункером (W. Brouncker, 1657), П. Ферма (Р. Fermat) и Дж. Валлисом (J. Wallis). Л. Эйлер (L. Euler) по недоразумению связал его с именем Дж. Пелля (J. Pell).





\end{document}